Let $ABC$ be a triangle with circumcircle $\Omega$.
The tangent to $\Omega$ at $A$ is called $t$.
Let $\omega_B$ be the circle outside of $ABC$, tangent to $t$ and to $AB$ at $B$.
Similarly, let $\omega_C$ be the circle outside of $ABC$, tangent to $t$ and to $AC$ at $C$.
Lastly, let $P$ be the intersection of the two common internal tangents
of $\omega_B$ and $\omega_C$, and let $\Gamma$ denote the circumcircle of the triangle $PBC$.

If the common tangents to $\Omega$ and $\Gamma$ intersect at a point $T$,
prove that $A$, $P$ and $T$ are collinear.

\emph{Remark:} A common internal tangent $\ell$ to two circles is tangent to both of them,
in such a way that the two circles \textbf{are not} on the same side of $\ell$.
